\documentclass{zhvt-classic}

% 使用本地源流明体（GenRyuMin2 TW）按文件路径加载，不依赖系统字体名
% 版框（内框／界栏／书口侧栏／鱼尾）改用与版心印章同一印泥棕色：
%   印章图片 yshs.png 实色 = RGB(135,68,52) = HTML 874434
\usepackage{xcolor}
\usepackage{tikz}
% sealbrown 双注册：xcolor 供 TikZ/pgf，L3 供 cls 版框
\definecolor{sealbrown}{HTML}{874434}
\ExplSyntaxOn
\color_set:nnn { sealbrown } { HTML } { 874434 }
\ExplSyntaxOff

\zhvtset{
  grid color       = sealbrown,
  main font        = GenRyuMin2TW-SB,
  main font path   = C:/Users/Administrator/texmf/fonts/opentype/public/GenRyuMin2TW/,
jiazhu font      = GenRyuMin2TW-M,
  jiazhu font path = C:/Users/Administrator/texmf/fonts/opentype/public/GenRyuMin2TW/,
  % 页码框容纳 4 字：正常页「九十七」+前置空 1 格=4 行恰好；
  % 100+ 页「一百零一」共 5 行超出，末字「一/二/三」被框裁掉（2026-10-04 实测），故放宽到 5
  shukou page number chars = 5,
  % 书口书名框高从默认 6 字减到 4 字，把省出的 2 字高拨给 chapter 框，
  % 使 6 字书名+卷名放进 9 字高的 chapter 框后不再压线（2026-10-05）
  shukou title chars = 4,
}

\title{史记}[张释之冯唐列传]
\author{汉·司马迁}
\maker{闲叟漫记·重製}


% ===== 索隐述赞样式（2026-11-08 四次定稿 · 愚大人亲手改定参数）=====
% ⚠ 竖排前提：自定义字体家族必须带 RawFeature={vertical:+vert}；
%   simhei.ttf 加 +vhal 会整体加载失败静默回退，只写 +vert；simfang 同理。
\newCJKfontfamily{\shuzanfont}[Path=C:/Windows/Fonts/,RawFeature={vertical:+vert}]{simhei.ttf}
\newCJKfontfamily{\shuzanbody}[Path=C:/Windows/Fonts/,RawFeature={vertical:+vert}]{simfang.ttf}
% 题头：152x30pt 圆角 6pt 棕底白字；baseline (0,2) 使整框落入界行格子不压版框线；
% 白字节点 x 不补偿、y +8.2pt（像素级墨迹实测 dx/dy≈0，L/R 留白约 20pt）。
\newcommand{\shuzanhead}[1]{%
  \begin{tikzpicture}[baseline={(0,2)}]
    \fill[sealbrown, rounded corners=6pt]
      (-76pt,-15pt) rectangle (76pt,15pt);
    \node[text=white, inner sep=0pt, anchor=center] at (0pt,8.2pt)
      {\shuzanfont\fontsize{23}{23}\selectfont #1};
  \end{tikzpicture}%
}
% ★ 修正书口卷名位置：chapter 框起始位置改回 cls 默认值，避免全局影响。
%   只在每个卷首页通过 \juan 命令里的空命令微调该页位置（安全，不会跨页污染）。
\ExplSyntaxOn
\AtBeginDocument{
  \dim_set:Nn \l__zhvt_shukou_chapter_title_start_dim
    { \l__zhvt_shukou_title_height_dim + \l__zhvt_font_size_dim * 2 }
  \dim_set:Nn \l__zhvt_shukou_page_number_start_dim
    { \l__zhvt_shukou_chapter_title_start_dim + \l__zhvt_shukou_chapter_title_height_dim }
}
\ExplSyntaxOff

\begin{document}

% ★★ 版心（书口）排布覆写：把书名放到鱼尾之下（2026-10-04）
%
%   cls 默认版心自上而下：书名（title 静态框，在鱼尾之上）→ 鱼尾 → 篇名（chapter 动态框）→ 页码
%   本模板改为：          鱼尾（位置不动）→ 书名 → 篇名 → 页码分割线 → 页码
%
%   手法：把 cls 放在 title 静态框里的书名，改放进 chapter 动态框（鱼尾之下）；
%         title 静态框留空。书名取 \l__zhvt_title_str，即 \title{...} 的第一参数，
%         所以换书时只改 \title，不必在此重复写书名。
%
%   ⚠ 三条必须遵守的写法（2026-10-04 实测）：
%     ① 汉字必须用 \hbox{...} 包住。flowfram 会把框内容装进一个宽 14.4mm（= 一个 line_sep
%        = 书口栏宽）的纵向 minipage 并按该宽度折行，而整页靠 PDF /Rotate 90 显示、
%        框内文字本身不变向 —— 不包 \hbox 时「围炉夜话」4 字会被逐字折成 4 行，
%        在版面上就成了横排 4 字并横向溢出版框。
%     ② 竖排下**一个 \hbox 就是一整列**。书名与篇名要同处一列，必须放进同一个 \hbox；
%        分成两个 \hbox 时，后一个会被当成新的一列、甩到隔壁栏去
%        （实测「序」跑到书口栏外侧 14.4mm 的正文列上）。中间的 \quad 在框内是横向空白，
%        显示出来正好是书名与篇名之间的纵向间隔。
%     ③ ★ 内容盒要写成 **与框等宽的定宽盒** `\hbox to \l__zhvt_shukou_chapter_title_height_dim{…\hfil}`。
%        实测：chapter 框是把内容盒沿纵向**居中**摆放的，内容宽度一变，书名就跟着漂：
%        自然宽度的 \hbox（书名 38.4mm）落在 x=157.2mm，加篇名后（57.6mm）落在 x=152.4mm，
%        目录页（无篇名）与正文页的书名就差半格。做成等宽盒后，无论有没有篇名，
%        内容固定从框顶起排，实测两页均落在 x=142.8mm，不再漂。
\ExplSyntaxOn
% 把书口 chapter 框（书名+篇名）起始位置下移 1 字高，避免 4 字卷名顶到鱼尾
% ⚠ cls 的 \__zhvt_calculate_dimensions: 在 \AtBeginDocument 里执行，会重算该 dim；
%   必须把覆写挂到 begin document 之后，否则会被二次覆盖（2026-10-05 实测）。
% (旧方案已移到导言区，直接修改 chapter 框起始位置)
% 卷标题命令 #1=书口篇名(鱼尾) #2=正文卷名（含卷序，卷名内空格用 ~）
%
% ★★ 目录页两项修订（2026-10-05，愚大人指出）★★
%   ① 页码起点统一对齐：
%      原状 —— 目录项写成 \par\hspace{4em}卷名\hspace{4em}页码，页码起点 = 卷名实宽 + 4em。
%      卷名实宽随字数变：「卷一~梁惠王上」6 字 +~、「卷七~离娄上」5 字 +~，
%      差 1 个字位 → 卷七~卷十 4 条的页码整体比其余 10 条高 27.21pt（一个字），看着上下不齐。
%      改法 —— 卷名补全角空格（宽 = 1 字 = 27.215pt）补足到 **7 字符位**再写进目录，
%      14 条页码起点遂全部落在同一条竖线上（起点对齐；末字随页码字数自然参差）。
%      ⚠ 只对目录里的副本 \l_tmpc_tl 补，正文卷名仍用 \l_tmpa_tl 原样 —— 正文像素不受影响。
%   ② 卷名整体上移半字位：前导 \hspace 由 4em 改 3.5em。
%      实测原状目录块在版框内偏下（天头留白 139.6pt / 地脚留白 110.7pt），
%      上移 0.5em = 13.6pt 后上下基本均分，视觉居中。
\NewDocumentCommand { \juan } { mm } {
	\clearpage
	\tl_set:Nn \l__zhvt_chapter_title_tl { #1 }
	\tl_set:Nn \l_tmpa_tl { #2 }
	\tl_set_eq:NN \l_tmpc_tl \l_tmpa_tl
	\int_step_inline:nn { 7 - \str_count:n { #2 } }
	  { \tl_put_right:Nn \l_tmpc_tl { 　 } }
	\tl_set:Nx \l_tmpb_tl { \zhnum { page } }
	\tl_put_left:Nn \l_tmpb_tl { \hspace { 4em } }
	\tl_put_left:Nn \l_tmpb_tl { \l_tmpc_tl }
	\tl_put_left:Nn \l_tmpb_tl { \par \hspace { 3.5em } }
	\addtocontents { toc } { \l_tmpb_tl }
	\skip_horizontal:n { \l__zhvt_font_size_dim * 4 }
	\l_tmpa_tl
	\par
	\thispagestyle { empty }
}
\cs_set_protected:Npn \__zhvt_set_shukou_content: {
	% ★ 鱼尾上方留白（title 静态框）嵌入闲叟钤印：
	%   整页 /Rotate 90，故此处把图用 \rotatebox{90} 转正为竖章；
	%   尺寸卡在版心条幅内（书内约 7mm 宽 × 33.8mm 高，留余量不压鱼尾）。
	\setstaticcontents * { title } {
		\vbox_to_ht:nn { \l__zhvt_baseline_skip_dim } {
			\vskip 14.2pt
			\hbox_to_wd:nn { \l__zhvt_shukou_title_height_dim } {
				\hfil
				\rotatebox [ origin = c ] { 90 } { \includegraphics [ width = 7mm , height = 33.8mm ] { yshs.png } }
				\hfil
			}
		}
	}
% ★ 版心（书口）统一改法（2026-11-16 愚大人定）：18pt＋框内居中＋书口文字同版框色
%   ① 原 27.2pt 时 7 字长篇名（如「廉颇蔺相如列传」）末字溢出 chapter 框、压住页码；
%      改 \fontsize{18pt}{18pt}，前后 \hfil 令字串在框内居中（不随篇名字数漂移）。
%   ② \vbox{\vskip 8.2pt} 定量下移，抵消缩字号后栏宽方向 8.2pt 上偏（像素实测 y 中心 543.3 = 原版）。
%      ⚠ \raisebox / \vspace* 在 flowfram 框内无效；换字号须同步换补偿：16pt→10.1、19pt→7.4。
%   ③ chapter 与 page~number 均置 \color{sealbrown}（HTML 874434），与版框／界行同色；正文仍黑。
%   ⚠ 页码字位、钤印、鱼尾、界行、正文一律未动；五部合订本同此口径。
	\setdynamiccontents * { chapter }
	{ \color{sealbrown} \vbox { \vskip 8.2pt \hbox to \l__zhvt_shukou_chapter_title_height_dim
		{ \hfil \fontsize{18pt}{18pt}\selectfont \l__zhvt_title_str \quad \l__zhvt_chapter_title_tl \hfil } } }
	\setdynamiccontents * { page~number } { \color{sealbrown} \quad \zhnum { page } }
}
\__zhvt_set_shukou_content:
\ExplSyntaxOff

\maketitle

% 单篇，无多卷目录，不渲染目录页

\juan{张释之冯唐列传}{卷一百零二~张释之冯唐列传第四十二}
\hspace*{2em}张廷尉释之者。堵阳人也。字季。有兄仲同居。以訾为骑郎。事孝文帝。十岁不得调。无所知名。释之曰。﹁久宦减仲之产。不遂。﹂欲自免归。中郎将袁盎知其贤。惜其去。乃请徙释之补谒者。释之既朝毕。因前言便宜事。文帝曰。﹁卑之。毋甚高论。令今可施行也。﹂于是释之言秦汉之间事。秦所以失而汉所以兴者久之。文帝称善。乃拜释之为谒者仆射。\\
\hspace*{2em}释之从行。登虎圈。上问上林尉诸禽兽簿。十馀问。尉左右视。尽不能对。虎圈啬夫从旁代尉对上所问禽兽簿甚悉。欲以观其能口对响应无穷者。文帝曰。﹁吏不当若是邪。尉无赖。﹂乃诏释之拜啬夫为上林令。释之久之前曰。﹁陛下以绛侯周勃何如人也。﹂上曰。﹁长者也。﹂又复问。﹁东阳侯张相如何如人也。﹂上复曰。﹁长者。﹂释之曰。﹁夫绛侯。东阳侯称为长者。此两人言事曾不能出口。岂斅此啬夫谍谍利口捷给哉。且秦以任刀笔之吏。吏争以亟疾苛察相高。然其敝徒文具耳。无恻隐之实。以故不闻其过。陵迟而至于二世。天下土崩。今陛下以啬夫口辩而超迁之。臣恐天下随风靡靡。争为口辩而无其实。且下之化上疾于景响。举错不可不审也。﹂文帝曰。﹁善。﹂乃止不拜啬夫。\\
\hspace*{2em}上就车。召释之参乘。徐行。问释之秦之敝。具以质言。至宫。上拜释之为公车令。\\
\hspace*{2em}顷之。太子与梁王共车入朝。不下司马门。于是释之追止太子。梁王无得入殿门。遂劾不下公门不敬。奏之。薄太后闻之。文帝免冠谢曰。﹁教兒子不谨。﹂薄太后乃使使承诏赦太子。梁王。然后得入。文帝由是奇释之。拜为中大夫。\\
\hspace*{2em}顷之。至中郎将。从行至霸陵。居北临厕。是时慎夫人从。上指示慎夫人新丰道。曰。﹁此走邯郸道也。﹂使慎夫人鼓瑟。上自倚瑟而歌。意惨凄悲怀。顾谓群臣曰。﹁嗟乎。以北山石为椁。用纻絮斫陈。蕠漆其间。岂可动哉。﹂左右皆曰。﹁善。﹂释之前进曰。﹁使其中有可欲者。虽锢南山犹有郄。使其中无可欲者。虽无石椁。又何戚焉。﹂文帝称善。其後拜释之为廷尉。\\
\hspace*{2em}顷之。上行出中渭桥。有一人从穚下走出。乘舆马惊。于是使骑捕。属之廷尉。释之治问。曰。﹁县人来。闻跸。匿桥下。久之。以为行已过。即出。见乘舆车骑。即走耳。﹂廷尉奏当。一人犯跸。当罚金。文帝怒曰。﹁此人亲惊吾马。吾马赖柔和。令他马。固不败伤我乎。而廷尉乃当之罚金。﹂释之曰。﹁法者天子所与天下公共也。今法如此而更重之。是法不信于民也。且方其时。上使立诛之则已。今既下廷尉。廷尉。天下之平也。一倾而天下用法皆为轻重。民安所措其手足。唯陛下察之。﹂良久。上曰。﹁廷尉当是也。﹂\\
\hspace*{2em}其後有人盗高庙坐前玉环。捕得。文帝怒。下廷尉治。释之案律盗宗庙服御物者为奏。奏当弃市。上大怒曰。﹁人之无道。乃盗先帝庙器。吾属廷尉者。欲致之族。而君以法奏之。非吾所以共承宗庙意也。﹂释之免冠顿首谢曰。﹁法如是足也。且罪等。然以逆顺为差。今盗宗庙器而族之。有如万分之一。假令愚民取长陵一抔土。陛下何以加其法乎。﹂久之。文帝与太后言之。乃许廷尉当。是时。中尉条侯周亚夫与梁相山都侯王恬开见释之持议平。乃结为亲友。张廷尉由此天下称之。\\
\hspace*{2em}後文帝崩。景帝立。释之恐。称病。欲免去。惧大诛至。欲见谢。则未知何如。用王生计。卒见谢。景帝不过也。\\
\hspace*{2em}王生者。善为黄老言。处士也。尝召居廷中。三公九卿尽会立。王生老人。曰﹁吾穇解﹂。顾谓张廷尉。﹁为我结韤。﹂释之跪而结之。既已。人或谓王生曰。﹁独柰何廷辱张廷尉。使跪结穇。﹂王生曰。﹁吾老且贱。自度终无益于张廷尉。张廷尉方今天下名臣。吾故聊辱廷尉。使跪结穇。欲以重之。﹂诸公闻之。贤王生而重张廷尉。\\
\hspace*{2em}张廷尉事景帝岁馀。为淮南王相。犹尚以前过也。久之。释之卒。其子曰张挚。字长公。官至大夫。免。以不能取容当世。故终身不仕。\\
\hspace*{2em}冯唐者。其大父赵人。父徙代。汉兴徙安陵。唐以孝著。为中郎署长。事文帝。文帝辇过。问唐曰。﹁父老何自为郎。家安在。﹂唐具以实对。文帝曰。﹁吾居代时。吾尚食监高袪数为我言赵将李齐之贤。战于钜鹿下。今吾每饭。意未尝不在钜鹿也。父知之乎。﹂唐对曰。﹁尚不如廉颇。李牧之为将也。﹂上曰。﹁何以。﹂唐曰。﹁臣大父在赵时。为官将。善李牧。臣父故为代相。善赵将李齐。知其为人也。﹂上既闻廉颇。李牧为人。良说。而搏髀曰。﹁嗟乎。吾独不得廉颇。李牧时为吾将。吾岂忧匈奴哉。﹂唐曰。﹁主臣。陛下虽得廉颇。李牧。弗能用也。﹂上怒。起入禁中。良久。召唐让曰。﹁公柰何众辱我。独无间处乎。﹂唐谢曰。﹁鄙人不知忌讳。﹂\\
\hspace*{2em}当是之时。匈奴新大入朝那。杀北地都尉卬。上以胡寇为意。乃卒复问唐曰。﹁公何以知吾不能用廉颇。李牧也。﹂唐对曰。﹁臣闻上古王者之遣将也。跪而推毂。曰阃以内者。寡人制之。阃以外者。将军制之。军功爵赏皆决于外。归而奏之。此非虚言也。臣大父言。李牧为赵将居边。军市之租皆自用飨士。赏赐决于外。不从中扰也。委任而责成功。故李牧乃得尽其智能。遣选车千三百乘。彀骑万三千。百金之士十万。是以北逐单于。破东胡。灭澹林。西抑彊秦。南支韩。魏。当是之时。赵几霸。其後会赵王迁立。其母倡也。王迁立。乃用郭开谗。卒诛李牧。令颜聚代之。是以兵破士北。为秦所禽灭。今臣窃闻魏尚为云中守。其军市租尽以飨士卒。私养钱。五日一椎牛。飨宾客军吏舍人。是以匈奴远避。不近云中之塞。虏曾一入。尚率车骑击之。所杀其众。夫士卒尽家人子。起田中从军。安知尺籍伍符。终日力战。斩首捕虏。上功莫府。一言不相应。文吏以法绳之。其赏不行而吏奉法必用。臣愚。以为陛下法太明。赏太轻。罚太重。且云中守魏尚坐上功首虏差六级。陛下下之吏。削其爵。罚作之。由此言之。陛下虽得廉颇。李牧。弗能用也。臣诚愚。触忌讳。死罪死罪。﹂文帝说。是日令冯唐持节赦魏尚。复以为云中守。而拜唐为车骑都尉。主中尉及郡国车士。\\
\hspace*{2em}七年。景帝立。以唐为楚相。免。武帝立。求贤良。举冯唐。唐时年九十馀。不能复为官。乃以唐子冯遂为郎。遂字王孙。亦奇士。与余善。\\
\hspace*{2em}太史公曰。张季之言长者。守法不阿意。冯公之论将率。有味哉。有味哉。语曰﹁不知其人。视其友﹂。二君之所称诵。可著廊庙。书曰﹁不偏不党。王道荡荡。不党不偏。王道便便﹂。张季。冯公近之矣。\\

% ===== 索隐述赞（司马贞《史记索隐》）· 统一样式 =====
% 题头「索隐述赞曰」五字并入反白框；正文每四句一列，定宽盒 \hfil 居列中；
% \fontsize{24}{41} 列距对齐全书界行（121+41k）；\clearpage 使述赞独占一页。
\clearpage
\shuzanhead{索隐述赞曰}\linebreak
{\shuzanbody\fontsize{24}{41}\selectfont
\hbox to \hsize{\hfil 张季未偶。见识袁盎。太子惧法。啬夫无状。\hfil}\linebreak
\hbox to \hsize{\hfil 惊马罚金。盗环悟上。冯公白首。味哉论将。\hfil}\linebreak
\hbox to \hsize{\hfil 因对李齐。收功魏尚。\hfil}%
\par}
\end{document}
